\chapter{Introduction}
\label{chapter:introduction}

\begin{introduction}
\begin{flushright}
``I'm not here to make you money. I'm here to make you rich.''
\end{flushright}

\vspace{0.5em}

This chapter motivates the study and defines its scope. It situates the work within the growing literature on Financial NLP, highlights the gap between sentiment-based return prediction and its integration into portfolio construction, and introduces sector-level ETFs as the chosen investment universe. The chapter closes with an outline of the thesis structure.
\end{introduction}

In today's financial landscape, vast amounts of unstructured and semi-structured textual data are generated and consumed at an unprecedented scale. Earnings calls, analyst reports, regulatory filings, macroeconomic commentary, financial news, and social media posts collectively constitute a rich and continuous stream of information that shapes market expectations and drives investor behavior. These diverse textual sources often carry nuanced signals that traditional numerical data alone cannot fully capture. As investors and financial institutions increasingly seek to make informed, real-time decisions, the ability to automatically interpret and act on this textual data has become not merely advantageous but essential~\cite{Loughran2016, Martina25}.

Natural Language Processing (NLP) has emerged as the primary toolkit for extracting structured knowledge from this flood of text, and its application to finance, a subfield increasingly referred to as Financial NLP (FinNLP), has expanded rapidly over the past decade. The use cases span a wide range of tasks and data sources: from annual reports and 10-K filings, where NLP methods extract signals about firm quality, management tone, and forward-looking uncertainty, to financial news articles processed to derive sentiment scores that predict short-term price movements, to social media platforms such as X, StockTwits, and Reddit, mined to capture the real-time pulse of retail investor sentiment \cite{3ac6b, 9d5c9, 2a15b, a12be, cc28c}. Across all these sources, a consistent empirical finding emerges: sentiment extracted from text provides complementary information to historical price data and technical indicators, improving predictive accuracy across equity markets, bond markets, and alternative assets alike \cite{2a15b, a12be, cc28c}.

Despite this body of evidence, the integration of NLP signals into portfolio optimization, as opposed to pure prediction tasks, remains comparatively underexplored. Portfolio optimization is concerned with the allocation of capital across an investable universe of financial assets to balance expected returns and risk, and while the field has evolved significantly over the decades, from classical mean-variance frameworks to deep learning and reinforcement learning approaches, the predominant inputs have remained mainly numerical: historical prices, returns, volatility estimates, and technical indicators. Systematically incorporating NLP-derived sentiment into the portfolio construction process itself, rather than merely into upstream return prediction, is a direction the literature is only beginning to explore \cite{Markowitz1952, 0df79, 3ac6b}.

\section{Research Objective and Scope}

The central objective of this thesis is to investigate whether NLP-derived features can enhance portfolio optimization at the sector level. The choice of sector-level allocation is deliberate. Individual stocks are too exposed to idiosyncratic noise (earnings surprises, management changes, litigation) that is orthogonal to the broader dynamics NLP signals are best equipped to capture. Broad market indices, by contrast, aggregate away the very cross-sectional variation that makes NLP-driven tilts actionable. Sectors occupy a more productive middle ground: large enough to attract extensive news coverage and online discussion, yet distinctive enough in their economic activities for sentiment signals to carry genuine sector-specific informational content. Both the interpretability of portfolio tilts and the quality of text-derived signals are, arguably, strongest at this level.

The empirical vehicle for testing these signals is the sector Exchange-Traded Fund (ETF). An ETF is a pooled investment vehicle that tracks an underlying index or basket of assets and trades continuously on an exchange, much like an individual stock \cite{hill2015etf}. Three properties make ETFs particularly well-suited to this research. First, they provide immediate within-sector diversification, attenuating the firm-specific volatility that would otherwise contaminate sector-level signals. Second, their intraday liquidity and low transaction costs make daily rebalancing practically feasible in a way that end-of-day mutual fund pricing would not permit. Third, and crucially for reproducibility, ETF holdings are publicly disclosed, so portfolio weights correspond unambiguously to known sector exposures. This transparency also supports the NLP pipeline at the heart of the study. Because each ETF maps cleanly to a defined sector of economic activity, textual data can be filtered and attributed to specific sectors in a principled and reproducible manner, allowing the use of sector-specific rather than market-wide sentiment features.

Against this backdrop, the thesis makes several novel contributions to the literature on NLP-driven asset allocation. It shifts empirical focus from the well-studied domains of individual stocks and broad indices to the underexplored sector level. It introduces sector ETFs as the research vehicle, exploiting their transparency, liquidity, and built-in diversification to construct a clean empirical setting. It develops a principled pipeline for attributing textual data to sectors, enabling the construction of meaningfully sector-specific features. Finally, it provides the first systematic empirical evaluation of whether such NLP-derived features improve sector-level portfolio performance, filling a clear gap in the literature and offering new evidence on the practical value of text-based signals in systematic investing.


\section{Thesis Structure}

The remainder of this thesis is organized as follows. Chapter~\ref{chapter:sota} reviews the theoretical and empirical foundations of the work, covering portfolio optimization paradigms, return prediction methods, and the literature on narrative analysis in finance, and concludes with the formal hypotheses. Chapter~\ref{chapter:materials} describes the data sources and collection procedures for ETF market data, financial news from GDELT, and social media content from Reddit. Chapter~\ref{chapter:methods} details the three portfolio optimization frameworks (a Deep Reinforcement Learning agent, a prediction-then-optimize pipeline, and an end-to-end Sharpe ratio maximizer), the technical and sentiment feature engineering pipeline, and the experimental design. Chapter~\ref{chapter:empirical_results} presents the empirical results for all frameworks and feature variants. Chapter~\ref{chapter:discussion} discusses the findings, benchmarks them against comparable work in the literature, reflects on their policy and practical implications, and identifies limitations and directions for future research. Final conclusions are drawn in Chapter~\ref{chapter:conclusion}.

% \bigskip
% \noindent\textit{All online resources cited in this thesis were last consulted in April 2026, unless otherwise indicated.}
